\chapter{MATERIALES Y MÉTODOS}

% ===========================================================================
% Tecnologías, arquitectura, diseño, desarrollo iterativo y plan de evaluación.
% Aquí van los diagramas (IT-46 secuencia scraping/chunking, IT-47 modelo de
% datos, IT-52 secuencia RAG, IT-53 storyboards) y el plan de validación
% científica definido ANTES de implementar.
% ===========================================================================

\section{Tecnologías y herramientas empleadas}

% Python, Scrapy, sentence-transformers, base vectorial, LLM en local, framework
% web. Cada elección relevante remite a su ADR.

\section{Arquitectura del sistema}

% IT-17. Diagrama del pipeline scraping -> limpieza -> chunking -> indexación ->
% recuperación -> generación. Explicar la independencia entre capas.

\section{Diseño de los datos}

% IT-47. Estructura de Asignatura, Chunk y resultado de recuperación. No hay un
% diagrama Entidad-Relación clásico porque no hay base de datos relacional;
% explicar por qué.

\section{Desarrollo: Iteración 0 (datos)}
% IT-18. Spider y chunker, referenciando issues y commits.

\section{Desarrollo: Iteración 1 (embeddings y vector DB)}
% IT-26.

\section{Desarrollo: Iteración 2 (pipeline RAG)}
% IT-33 (parte de desarrollo). Incluir el diagrama de secuencia del flujo RAG (IT-52).

\section{Desarrollo: Iteración 3 (aplicación web)}
% IT-37. Incluir storyboards (IT-53).

\section{Plan de evaluación científica}

% Definido ANTES de implementar. Métricas de recuperación (Recall@K, MRR) y de
% generación (faithfulness, answer relevancy). Conjunto de preguntas etiquetadas
% (IT-20) como verdad de referencia. Análisis de amenazas a la validez (IT-32).
